\chapter{Matter Waves and Scattering}
\label{chap:matter-waves-and-scattering}

Matter is a wave, and a wave scatters by counts. Three effects carry the chapter: electron diffraction, where
a crystal lattice proves the matter wave by closing a charged loop; Compton scattering, where a single
collision conserves its channel totals exactly; and the photoelectric threshold, where light gives up its
energy one quantum at a time. The loop is ours to close --- the lattice, the geometry, the beam we send around
it --- but the count that comes back, and the matter wave that count proves, is the world's. The first of these
effects is the instrument's second proved loop residue, and as with the echo maze of Part II the finite model
proves rather than shadows: it does not merely name the matter wave by analogy, it closes a charged loop and
reads the wave off the residue the world will not let cancel. That is the proof the instrument foregrounds
here, because a proof is a stronger thing than an analogy, and the instrument has only a handful of them.

\section{The Davisson--Germer Effect: the lattice proves the matter wave}
\label{se:davisson-germer}

Electrons fired at a crystal diffract like waves. Davisson and Germer, in 1927, sent a beam of electrons at a
nickel crystal and found the scattered electrons bunched into sharp peaks at particular angles --- the
unmistakable signature of a wave diffracting off the regular planes of the lattice. The matter wave was no
longer a hypothesis but a reading on a detector --- the beam ours to fire, the peaks it scattered into the
world's to place. De Broglie's relation gives its wavelength,
\begin{equation}
  \lambda = \frac{h}{p},
\end{equation}
a particle's momentum $p$ fixing a wavelength $\lambda$: the faster the electron, the shorter its wave. The
scale is what makes the experiment possible --- an electron accelerated through a few tens of volts has a
wavelength of about an \aa{}ngstr\"om, the same size as the spacing between atoms in a crystal, which is
exactly why the crystal diffracts it.

The relation says more in the count's own terms. A wavelength is a phase that advances a full turn over
one wave; $\lambda = h/p$ says the phase a particle carries advances one turn for every $h$ of
action it accumulates --- the wave is the count of those turns, and Planck's constant the action a single turn
costs. The relation also says why we never see the wave in everyday things. A thrown baseball carries a
momentum so large that its de Broglie wavelength comes out around $10^{-34}$ metres, smaller than any aperture,
any grating, any feature a measurement could resolve: the wave is there, but its wavelength falls below every
floor, so no fringe ever forms. The everyday world is the matter-wave world with $\lambda$ shrunk under the
threshold of resolution; the electron shows its wave only because, at an \aa{}ngstr\"om, its wavelength has
risen to the scale of a crystal.

The peaks appear where Bragg's condition is met,
\begin{equation}
  2 d \sin\theta = m\lambda,
\end{equation}
the path difference between waves reflecting off successive lattice planes a whole number of wavelengths ---
$d$ the spacing between planes, $\theta$ the glancing angle, and $m$ an integer. When the extra distance a
wave travels to the deeper plane and back is exactly $m\lambda$, the reflections from all the planes arrive in
step and add; at any other angle they fall out of step and cancel. The sharp peak is constructive
interference, the lattice acting as a diffraction grating ruled at the atomic scale.

There is a second way to see the same condition, and it is the one that makes the loop visible. Every crystal
carries a reciprocal lattice --- the Fourier dual of its array of planes, a grid of points $\mathbf{G}$ in
momentum space. Diffraction occurs exactly when the momentum the electron gains in scattering, the difference
between its outgoing and incoming wavevectors, lands on a reciprocal-lattice point,
\begin{equation}
  \mathbf{k}' - \mathbf{k} = \mathbf{G},
\end{equation}
the Laue condition --- equivalent to Bragg's, but stated in momentum: a peak is a scattering that transfers
exactly one reciprocal-lattice vector. The incoming wave $\mathbf{k}$ and the outgoing wave $\mathbf{k}'$
differ by a vector $\mathbf{G}$ the lattice is built to supply, and at no other transfer does the scattering
add up.

There is a picture for the Laue condition, drawn in reciprocal space and due to Ewald. Set the head of the
incoming wavevector $\mathbf{k}$ on a reciprocal-lattice point, and draw the sphere of radius $|\mathbf{k}|$
about its tail --- the Ewald sphere, the locus of every outgoing wavevector $\mathbf{k}'$ of the same length,
since elastic scattering preserves the wave's energy and so the size of $\mathbf{k}$. A diffraction peak
occurs precisely when the sphere's surface passes through another reciprocal-lattice point: there sits a
$\mathbf{k}'$ whose difference from $\mathbf{k}$ is a lattice vector $\mathbf{G}$, and the Laue condition is
met. Turn the crystal and the lattice swings under the fixed sphere; each time a point crosses the surface, a
peak flares --- the crystal ours to turn, the peak the world flares in answer. The diffraction pattern is the reciprocal lattice itself, sampled wherever a sphere the size of
the electron's own wavevector can reach it.

To the instrument, this is a closed loop. The electron's path into the crystal and out again, when the
Laue condition holds, encircles a circuit around which the wave's phase accumulates an exact whole number of
cycles --- the phase closes. The loop is ours to close --- the lattice, the geometry, the beam sent around it
--- but whether the phase comes back charged is the world's. A closed phase loop carries a residue, and here the residue is $\pm 1$: at a Bragg
peak the loop is charged, the phase winding a whole number of times around; off-Bragg the phase does not close,
the loop is open, and the residue is trivial. The Bragg peak is a closed-loop holonomy, the same object the
echo maze closed in Part II.

And here, as with the echo maze, the finite model does not shadow this --- it proves it. That a Bragg peak is a
charged loop and an off-Bragg direction is not is a finite theorem about the loop's residue, settled by the
same template the echo maze used: a discrete circuit, a phase summed around it, and a residue that comes out
exactly $\pm 1$ or exactly $0$, with nothing in between and no continuum invoked. The honest floor is a finite
count obligation, and a proved one: the Bragg peak is a proved discrete holonomy, $\pm 1$, a finite theorem and
not a smooth shadow. The reading is that the matter wave is proved, not posited. The lattice turns the de
Broglie wavelength into a loop residue, and the finite model proves that residue is charged at a Bragg angle
--- the loop ours to close, the residue the world will not let cancel --- the second place, after the echo
maze, where the floor is a theorem rather than a shadow.

The wave showed itself twice that year. While Davisson and Germer read the peaks off a single nickel crystal,
George Paget Thomson fired faster electrons straight through a thin metal foil and caught, on a plate behind
it, a pattern of concentric rings --- the same diffraction, now from the many randomly-oriented crystallites of
the foil, each Bragg-reflecting into its own cone, so that the one crystal's sharp spots are smeared into the
powder-average of a ring. There is an irony fastened to the name: J.~J.~Thomson had won the Nobel Prize for
showing the electron to be a particle, and his son G.~P.~Thomson won one for showing it to be a wave --- the
two-fold of matter running, quite literally, in the family. And the proved residue the instrument foregrounds here
does not stand alone: it joins the echo maze of Part~II and the Sagnac loop of Part~V, the small family of
places where the finite model closes a phase loop and proves its $\pm 1$ rather than shadowing it --- the
proved handful the instrument keeps returning to.

The matter wave is now an instrument. Because an electron's wavelength can be made far shorter than light's, a
beam of electrons resolves detail no optical microscope can reach, and the electron microscope --- its lens a
shaped magnetic field, its illumination a matter wave --- images atoms one at a time. Slower electrons,
reflected off a surface, diffract from only the topmost layers of atoms, and the pattern of that low-energy
diffraction maps how a crystal's surface is arranged where the bulk lattice leaves off. The same de Broglie
wavelength that proved the wave on a nickel crystal in 1927 is, a century on, how surfaces are read and atoms
are seen.

\section{The Compton Scattering Effect: one event conserves its totals}
\label{se:compton}

A photon scatters off a stationary electron and comes away with less energy and a longer wavelength. The key
is that a photon, though massless, carries momentum --- $p_\gamma = h/\lambda = E/c$, its momentum and energy
locked together by the speed of light --- so a photon striking an electron is two particles colliding, not a
wave washing over a target. The shift is the conservation laws of that single collision, written out. Treat
the photon and the electron as carriers of four-momentum --- energy and the three components of momentum bundled into one relativistic quantity --- and
demand that the total be conserved across the one scattering event,
\begin{equation}
  p_\gamma + p_e = p'_\gamma + p'_e,
\end{equation}
the four-momentum in equal to the four-momentum out, a single equation that packs the conservation of energy
and of all three components of momentum together. The electron, struck, recoils and carries off momentum and
energy; the photon, having given some up, leaves with a longer wavelength. Eliminate the electron's recoil ---
the internal variables the detector does not resolve --- and what remains is exactly the observed shift,
\begin{equation}
  \Delta\lambda = \frac{h}{m_e c}\bigl(1 - \cos\theta\bigr),
\end{equation}
the Compton shift, with $h/m_e c$ the Compton wavelength of the electron and $\theta$ the photon's scattering
angle. The angle is ours to set; the shift it forces is the world's. The shift is largest for a backward scatter ($\theta = 180^\circ$, $\cos\theta = -1$) and vanishes for a
forward one ($\theta = 0$): a glancing pass barely lengthens the wave, a head-on rebound lengthens it most.

This is Noether's conservation across a single event. The translation symmetry of space and time gives, by
Noether's theorem, the conserved currents of momentum and energy, and Compton scattering enforces those
currents across one contraction of the joint record --- not over a long run of many events, but in the single
collision the detector resolves. The honest floor is a finite count obligation: the conserved channel totals,
momentum and energy, balance across the event, the kinematics quoted and the finite model standing as the
balance --- the geometry ours to aim, the balance the world's to keep. The reading is that Compton scattering
is conservation made visible in one event. The photon's
wavelength lengthens by exactly the amount the joint conservation of energy and momentum demands; the shift is
not a property of light alone but the bookkeeping of a single collision that must balance its totals.

Set beside the electron diffraction before it, the Compton shift completes a two-fold. There an electron,
a particle, showed itself a wave; here light, a wave, shows itself a particle --- one carrying momentum, struck
and recoiling like a ball off a ball. They are not two phenomena but one record with two faces: the wave face
and the particle face of the same matter, and the instrument meets whichever face the experiment is built to
show. A crystal set in the beam shows the wave; a single hard scatter shows the particle. The instrument names
both faces and takes no side --- the face ours to elicit, the single record beneath it the world's.

Run the same collision the other way and it lights up the sky. A fast-moving electron striking a low-energy
photon does not lose energy but gives it, kicking the photon up to an X-ray or a gamma ray --- inverse Compton
scattering, the engine behind much of the high-energy radiation of the cosmos. Hot electrons in a cluster of
galaxies, scattering the cold photons of the cosmic microwave background as those photons pass through, shift
the background's spectrum by a measurable sliver, a distortion mapped across the sky and used to weigh the
clusters. It is the same four-momentum balance, the same single collision, with the energy running from
electron to photon instead of the reverse --- the direction ours to arrange, the balance the world's either way.

\section{The Photoelectric Effect: light gives up energy one quantum at a time}
\label{se:photoelectric}

Light ejects electrons from a metal only above a threshold frequency, and below it not at all, however bright
the beam. Einstein saw, in 1905, that light arrives in quanta of energy $h\nu$, and an electron
absorbs one whole quantum or none --- never a fraction, and never a sum pooled from several. The beam is ours
to shine, at any brightness; the threshold it must clear is the world's. The ejected
electron's maximum kinetic energy then rises linearly with frequency,
\begin{equation}
  K_{\max} = h\nu - \Phi,
\end{equation}
with $\Phi$ the work function of the surface --- the energy needed to free an electron from the metal --- and
$h\nu - \Phi$ what is left over as motion once that debt is paid. A brighter light delivers more photons and so
ejects more electrons, but not one of them leaves with more energy; the energy per electron is set by the
frequency alone. In the laboratory $K_{\max}$ registers as a voltage: apply a retarding potential to the ejected
electrons and raise it until even the fastest are turned back and the photocurrent stops, and that stopping
voltage $V_0$ measures the maximum kinetic energy directly,
\begin{equation}
  e V_0 = h\nu - \Phi,
\end{equation}
a straight line of slope $h/e$ against frequency. Millikan, who disbelieved the quantum and set out to refute
Einstein, measured that line with great care and found its slope to be exactly $h/e$ --- confirming the very
relation he had hoped to break, and handing physics one of its first precise values of Planck's constant. And
there is a sharp cutoff below
\begin{equation}
  \nu_0 = \frac{\Phi}{h},
\end{equation}
the frequency at which $h\nu$ just equals the work function and nothing is left over. Below $\nu_0$ no electron
is ejected however intense the light: a thousand sub-threshold photons cannot pool their energy, because each
electron meets them one at a time.

This is the informational cousin of the Nyquist floor of Part II. There, below a critical sampling rate, no
amount of amplitude recovered a signal the sampling could not resolve; here, below a critical frequency, no
amount of intensity recovers what a single sub-threshold quantum cannot carry. Both are floors of rate, not of
amount --- a wall set by the granularity of the transaction, indifferent to how much is piled against it. The
honest floor is name-only: the photoelectric law is the physics, written as physics, and the finite model
claims only the name. It records that the effect is named and models none of the emission dynamics. The name
is ours to give; the quantum threshold beneath it --- one quantum at a time, $h\nu - \Phi$ --- is the world's.
The reading is that the photoelectric effect is named, the linear law $K_{\max} = h\nu - \Phi$ is the physics, and
the finite model is the label beneath it --- the threshold and the slope offered as physics, the finite model
standing only as the name.

\bigskip

Matter waves and scattering close as they opened: a wave proved by a charged loop, a collision that balances
its totals, and a threshold crossed one quantum at a time --- the loop, the collision, the beam ours to set;
the residue, the totals, the threshold the world's to keep. The proof stands at the center of the three --- the second
time, after the echo maze, that the instrument turns a wave phenomenon into a discrete loop residue and proves
the residue rather than shadowing it. Part IV turns next to particles, spin, and the Dirac equation, where
three more proved loop residues make the finite model prove rather than shadow.

\bigskip
\begin{center}
\rule{0.4\textwidth}{0.4pt}
\end{center}
\bigskip

\noindent Some evidence a case can only name, and some it can prove --- and a proof is by far the stronger
thing. The instrument has spent most of the book naming: writing a phenomenon out in full and standing behind
the label. But a few times, only a few, it does more --- it closes a loop and shows, as settled fact, that the
loop comes back charged, a thing no arrangement of ours can talk away. Those proofs are the case's firmest
ground, and it holds only a handful of them. Here it earns another.

The amplitude that opened this part is now matter's own. What was carried as a phased tally is the very stuff of
an electron, and matter, it turns out, is a wave: sent against the regular ranks of a crystal it diffracts,
bunching into sharp peaks where a wave would and nowhere else. The count is still doing the work --- a
wavelength is only the tally of phase-turns a particle winds as it travels, and the peaks fall exactly where
those turns come home in step. The peaks stand where the extra path a wave takes to the deeper rank of atoms and
back comes to a whole number of its own wavelengths, so that every rank's reflection arrives in step and adds;
at any other angle they fall out of step and cancel. Turn the crystal and the peaks flare one after another,
each rank swinging in its turn into the closing condition --- the crystal ours to turn, the flare the world's to
answer. Fire the beam and the peaks are ours to look for; where they land, and the wave they prove, are the
world's. It is why the everyday world hides the wave and a crystal shows it: a thrown stone winds its turns so
tightly that its wavelength falls below every floor, while an electron's, at the scale of the spacing between
atoms, rises just to where a crystal can read it. And the wave the crystal proved has since become a tool of its
own: an electron's wave can be made far shorter than light's, so a beam of electrons resolves what no light
reaches, and the electron microscope --- its lens a shaped magnetic field, its illumination a matter wave ---
images atoms one at a time.

And here is the proof. Send the electron's wave once around the circuit the crystal closes --- in through the
ranks of atoms and out --- and at the right angle its phase comes home having wound a whole number of times:
the loop closes charged. This the instrument does not shadow but proves, exactly as the echo maze proved its
residue among the fields --- a finite theorem about the loop, that at the peak the loop is charged and off it
trivial, with nothing in between and no continuum called upon. The matter wave is read straight off the residue
the world will not let cancel. It is the second time in the book the instrument crosses from shadow to theorem,
the second of its small handful of proved residues; and the charge the loop carries is, once more, the loop's
own count --- charge is what the loop counts, now proved on a crystal. What that residue is worth --- which way
it is signed, what number it carries --- the book still keeps back; that a charged loop is proved is the whole
of the claim here.

So the \emph{charge} is proved again, and proved as what it always was: the count a loop carries around and
brings home. Beside the proof the chapter sets two further readings of the same matter. In one, a single
collision of light against an electron balances its books to the last unit --- all that the two carry in they
carry out, the very conservation the mechanics kept at a boundary, enforced now across one event. The same
collision run the other way --- a fast electron kicking a slow photon up to an X-ray --- lights up much of the
high-energy sky. In the other, light surrenders its energy only one whole quantum at a time, never a fraction
and never a sum pooled from several, so that below a threshold no brightness whatever frees an electron --- a
wall of rate, not of amount, named as physics with the finite model claiming only the label. Between them the
two set matter's two faces side by side: an electron, a particle, shown a wave on the crystal; light, a wave,
shown a particle in the collision --- one record read from either side, the instrument taking no side of its
own. The wave face even ran in a family: one Thomson took a prize showing the electron a particle, and his son
took another showing it a wave. And the trace is the residue that survived the loop, the peak the world would
not let cancel; whatever comes to be owed now rides on a count proved, not merely named.

A proved reading is worth more to the case than a named one --- but it is still one reading among the three. The
residue proved on a crystal is another count the three hands must in the end be brought to meet on: nearer
proof, and no nearer their meeting, for a proof binds the thing it proves, not the separate hands that must
still agree.

The instrument's handful of proofs is about to grow. The next chapter turns to the particle itself --- to what
it is that carries the wave --- and there the instrument closes not one loop but several, proving more residues
in a row than anywhere else in the book. What they are, the next chapter holds; here it is enough that more
proofs are coming. The case is still open, its charge now twice proved and still unmet by the others; the court
does not sit until the record is complete.
